\documentclass[10pt,a4paper]{amsart}
\usepackage[margin=19mm]{geometry}
\usepackage{amsmath,amssymb,mathtools}
\usepackage[scr=rsfs,cal=euler]{mathalfa}
\usepackage{xfrac}
\usepackage[unicode,hidelinks,bookmarks=false]{hyperref}
\newcommand{\ie}{i.e.\ }
\newcommand{\cf}{cf.\ }
\newcommand{\resp}{respectively\ }
\newcommand{\Subsection}{\subsection}
\providecommand{\nobreakdash}{\nobreak\hbox{-}}
\setlength{\emergencystretch}{2em}
\multlinegap0pt
\usepackage{amssymb}
\usepackage[shortlabels]{enumitem}
\usepackage{mathtools}
\usepackage{cancel}
\usepackage{tikz}
\newtheorem{theorem}{Theorem}
\title{Infinite-dimensional nonholonomic and vakonomic systems}
\author{Alexander G. Abanov, Boris Khesin}
\date{Source: arXiv:2511.00629v2 (CC BY 4.0)}
\begin{document}
\maketitle

\noindent\emph{Source and licence.} Infinite-dimensional nonholonomic and vakonomic systems. Version arXiv:2511.00629v2 (CC BY 4.0), distributed by arXiv under the Creative Commons Attribution 4.0 International licence, http://creativecommons.org/licenses/by/4.0/, as recorded in the arXiv licence metadata for this exact version and shown on the abstract page https://arxiv.org/abs/2511.00629v2. Full licence notice: \texttt{licenses/arxiv-2511.00629-CC-BY-4.0.txt}.\par\smallskip

\noindent\textit{Editorial scope.} Counted unit: the article's theorem theorem (source0.tex lines 1350--1353). The article's complete original proof of it is delivered with it (source0.tex lines 1355--1362, one \texttt{proof} environment of 8 lines). No mathematics is written, repaired or re-derived: the statement and the proof are the article's own lines, in the article's own notation, and no retained line is substituted or re-flowed.  No further source-local context is needed: every cross-reference of every retained line resolves inside the two delivered spans.  Nothing was omitted from any retained span, so the package carries no omission record.  No retained line contains a citation, so the wrapper carries no bibliography and no bibliography entry.  No retained line contains a figure environment, an image-inclusion command or drawing code, so no image is delivered and the package declares no asset dependency.  Editorial formatting only, in addition: an AMS \texttt{amsart} preamble that restates the article's own declarations and macros used by the retained lines, namely the environments \texttt{theorem} on separate counters, together with the standard packages those lines need.  The retained environments print in this document as Theorem 1 in order of appearance, whereas the article prints the same items with the numbering of its own full document.  The complete native source of the article as distributed in the arXiv e-print for this exact version is bundled unchanged as \texttt{sources/188263/source0.tex}.
\begin{theorem}
The snake-type motion with the  infinite-dimensional skate constraint
corresponds to a curve subordinated to an  infinite-dimensional  Goursat structure, the Cartan distribution in the infinite jet space.
\end{theorem}

\begin{proof}
 Indeed, consider the curve $z(s, t_*)$ at any moment $t_*\in (0,1)$. At the moment $t_*+\delta t$ the curve
$z(s, t_*+\delta t)$ has the same prolongation in $s$ at the curve $z(s, t_*)$, since all its derivatives are predefined
by the one-dimensional image in $\mathbb R^2$. Hence the motion in $t$ can be regarded as the motion
along its prolongation of the fixed curve in the plane, while the prolongation must be everywhere tangent to the Cartan distribution in the infinite jet space. Furthermore, a curve in the plane can be locally regarded as
the graph of a function $\mathbb R\to \mathbb R$. 
Thus the statement  follows from one-dimensionality of the image in the plane and the properties of the Cartan distribution for functions of one variable.   
\end{proof}

\end{document}
